\begin{proposition}[5.13.3]
\label{IV.2.s5.13.3-fr}
\begin{enumerate}
\item[\textup{(i)}] Para todo índice $\alpha$, sea $\mathfrak p_\alpha$ un ideal
primo de $A_\alpha$ tal que, para $\alpha\leq\beta$, se tenga
$\varphi_{\beta\alpha}(\mathfrak p_\alpha)\subset\mathfrak p_\beta$.
Entonces $\mathfrak p=\varinjlim\mathfrak p_\alpha$ es un ideal primo de $A$.
En particular, si los $A_\alpha$ son íntegros, lo mismo ocurre con $A$.
\item[\textup{(ii)}] Si, además, se tiene
$\mathfrak p_\alpha=\varphi_{\beta\alpha}^{-1}(\mathfrak p_\beta)$ para
$\alpha\leq\beta$, entonces $A_{\mathfrak p}$ se identifica canónicamente
con $\varinjlim(A_\alpha)_{\mathfrak p_\alpha}$.
\item[\textup{(iii)}] Si $A$ es un anillo local, $\mathfrak p$ su ideal maximal, y si
los homomorfismos $\varphi_\alpha$ son inyectivos, entonces, poniendo
$\mathfrak p_\alpha=\varphi_\alpha^{-1}(\mathfrak p)$, se tiene
$A=\varinjlim(A_\alpha)_{\mathfrak p_\alpha}$ y los homomorfismos
$(A_\alpha)_{\mathfrak p_\alpha}\to A$ son inyectivos.
\end{enumerate}
\end{proposition}

Las dos afirmaciones de (i) son de hecho equivalentes, puesto que
$A/\mathfrak p=\varinjlim(A_\alpha/\mathfrak p_\alpha)$. Supongamos, pues,
que todos los $A_\alpha$ son íntegros, y sean $x$, $y$ dos elementos de $A$
tales que $xy=0$. Existe entonces un índice $\alpha$ tal que
$x=\varphi_\alpha(x_\alpha)$, $y=\varphi_\alpha(y_\alpha)$, con
$x_\alpha,y_\alpha$ en $A_\alpha$, y
$\varphi_\alpha(x_\alpha y_\alpha)=0$; pero esto implica la existencia de un
índice $\beta\geq\alpha$ tal que, si se pone
$x_\beta=\varphi_{\beta\alpha}(x_\alpha)$,
$y_\beta=\varphi_{\beta\alpha}(y_\alpha)$, se tenga $x_\beta y_\beta=0$;
como $A_\beta$ es íntegro, se tiene $x_\beta=0$ o $y_\beta=0$ y, como
$x=\varphi_\beta(x_\beta)$ e $y=\varphi_\beta(y_\beta)$, esto termina de
demostrar que $A$ es íntegro.

Demostremos ahora (ii). La hipótesis implica que los
$A_\alpha\setmin\mathfrak p_\alpha$ forman un sistema inductivo de partes
multiplicativas de los $A_\alpha$, cuyo límite inductivo es
$A\setmin\mathfrak p$; por consiguiente, los anillos locales
$(A_\alpha)_{\mathfrak p_\alpha}$ forman un sistema inductivo (en el que los
homomorfismos de transición son \emph{locales}), y los homomorfismos canónicos
$(A_\alpha)_{\mathfrak p_\alpha}\to A_{\mathfrak p}$ un sistema inductivo de
homomorfismos, cuyo límite inductivo
$\psi:\varinjlim((A_\alpha)_{\mathfrak p_\alpha})
\longrightarrow A_{\mathfrak p}$
es un homomorfismo sobreyectivo. Queda por ver que $\psi$ es inyectivo. Ahora
bien, si $x_\alpha\in A_\alpha$ y
$s_\alpha\in A_\alpha\setmin\mathfrak p_\alpha$ son tales que
$\varphi_\alpha(x_\alpha)/\varphi_\alpha(s_\alpha)=0$ en
$A_{\mathfrak p}$, existe $s'\in A\setmin\mathfrak p$ tal que
$s'\varphi_\alpha(x_\alpha)=0$, y, reemplazando si es preciso $\alpha$ por un
índice mayor, se puede suponer que $s'=\varphi_\alpha(s'_\alpha)$, con
$s'_\alpha\in A_\alpha\setmin\mathfrak p_\alpha$, de donde
$\varphi_\alpha(s'_\alpha x_\alpha)=0$; existe, por tanto, un
$\beta\geq\alpha$ tal que
$\varphi_{\beta\alpha}(s'_\alpha x_\alpha)=0$ en $A_\beta$, lo que implica
$\varphi_{\beta\alpha}(x_\alpha)/\varphi_{\beta\alpha}(s_\alpha)=0$ en
$(A_\beta)_{\mathfrak p_\beta}$, y prueba, por tanto, que la imagen canónica de
$x_\alpha/s_\alpha$ en $\varinjlim((A_\alpha)_{\mathfrak p_\alpha})$ es nula.

Finalmente, la primera afirmación de (iii) resulta de (ii) y de que
$A_{\mathfrak p}=A$. Además, los elementos de
$A_\alpha\setmin\mathfrak p_\alpha$ son invertibles en $A$, por lo que también se
tiene $A_{\mathfrak p_\alpha}=A$, y el hecho de que
$(A_\alpha)_{\mathfrak p_\alpha}\to A$ sea inyectivo se deduce, por planitud,
de la hipótesis de que $\varphi_\alpha$ es inyectivo.

\begin{corollary}[5.13.4]
\label{IV.2.s5.13.4-fr}
Supongamos que todos los $A_\alpha$ son íntegros, y que los
$\varphi_{\beta\alpha}$ son inyectivos. Entonces, si $A'_\alpha$ es la
clausura integral de $A_\alpha$, $A'=\varinjlim A'_\alpha$ es la clausura
integral de $A$. En particular, si $A_\alpha$ es integralmente cerrado para
todo $\alpha$, lo mismo ocurre con $A$. Si los $A_\alpha$ son anillos locales
unirramosos \textup{(resp.} geométricamente unirramosos) \textup{(\textbf 0, 23.2.1)}, y los
$\varphi_{\beta\alpha}$ homomorfismos locales inyectivos, $A$ es unirramoso
\textup{(resp.} geométricamente unirramoso).
\end{corollary}

Si $K_\alpha$ es el cuerpo de fracciones de $A_\alpha$, resulta de (5.13.3,
(ii)), aplicado a los $\mathfrak p_\alpha=0$, que
$K=\varinjlim K_\alpha$ es el cuerpo de fracciones de $A$, y
$A'=\varinjlim A'_\alpha$ se identifica con

\oldpage[IV-2]{133}

un subanillo de $K$, evidentemente entero sobre $A$. Recíprocamente, sea
$z\in K$ un elemento que verifica una ecuación de dependencia integral
$z^n+\sum_{i=1}^{n}c_i z^{n-i}=0$,
donde los $c_i\in A$. Existe un índice $\alpha$, un
$z_\alpha\in K_\alpha$ y unos $c_{i\alpha}\in A_\alpha$ tales que $z$
(resp. los $c_i$) sean las imágenes de $z_\alpha$ (resp. de los
$c_{i\alpha}$), por tanto se tiene $z_\alpha\in A'_\alpha$, lo que prueba que
$z\in A'$. Se deduce inmediatamente que, si los $A_\alpha$ son integralmente
cerrados, lo mismo ocurre con $A$. Si los $A_\alpha$ son locales y los
$\varphi_{\beta\alpha}$ homomorfismos locales, se sabe
$(\mathbf{0}_{\mathrm{III}}, 10.3.1.3)$ que $A$ es un anillo local y que, si
$k_\alpha$ es el cuerpo residual de $A_\alpha$,
$k=\varinjlim k_\alpha$ es el de $A$; si los $A'_\alpha$ son locales, $A'$ es
local y, por tanto, $A$ es unirramoso; si además el cuerpo residual
$k'_\alpha$ de $A'_\alpha$ es una extensión radicial de $k_\alpha$ para todo
$\alpha$, $k'=\varinjlim k'_\alpha$, que es el cuerpo residual de $A'$, es
una extensión radicial de $k$, y, por tanto, $A$ es geométricamente
unirramoso.
